\subsection{一个形式做法及其用途}\label{uses}

本节讨论一个形式化的量子化配方。它并不回避反常（反常确实无法回避），但仍为此后的论证提供了有用的方向。

给定经典相空间 $M$ 上的预量子线丛 $\frL$，其联络为 $\sT/\hbar$，曲率为 $\omega=\d \sT /\hbar$，则可以写出一个经典作用量，其相应的路径积分将在形式上解决 $M$ 的量子化问题。对 $M$ 中两点 $r,r'\in M$ 之间的路径 $\gamma\subset M$，作用量就是
\be\label{wiff}\sI=\frac{1}{\hbar}\int_\gamma \sT .\ee
或者，令 $I=[0,1]$ 为单位区间，把 $\gamma$ 看成映射 $x:I\to M$ 的像，则作用量可写成
\be\label{iff} \sI =\frac{1}{\hbar}\int_I  x^*(\sT ). \ee
注意 $\sI$ 在区间 $I$ 的重参数化下不变。

若 $I$ 的端点 $0,1$ 映到点 $r,r'\in M$，则为了把作用量 $\sI$ 定义为一个数，需要在这些点处平凡化 $\frL$ 的纤维。把这些纤维记作 $\frL_r$、$\frL_{r'}$。设 $\pP_{r,r'}$ 是端点映到 $r,r'$ 的映射 $x:I\to M$ 的空间。考虑 Feynman 积分
\be\label{hiff} Z(r,r')=\int_{\pP_{r,r'}}\D x \,\exp(\i \sI). \ee
对固定的 $x$（在试图对 $x$ 积分之前），被积函数 $e^{\i \sI}$ 是从 $\frL_r$ 到 $\frL_{r'}$ 的映射；人们可以期望对 $x$ 的积分给出映射 $Z_{r,r'}:\frL_r\to \frL_{r'}$，并由它构造出量子化相空间 $M$ 的 Hilbert 空间。

按字面看，这不能真的成立，因为这里 $r$ 与 $r'$ 将对应初始态与末态，而初始与末态量子态只应依赖相空间变量的一半。没有更多输入，不应能构造出与 $M$ 中任意两点相联系的映射 $Z_{r,r'}$。事实上，经典力学的对称群 $\G$（即第 \ref{problem} 节所述的、带 $\frL$ 上作用提升的恰当辛同构）如此之大，以至于不可能有 $\G$-不变的 $Z_{r,r'}$ 定义。群 $\G$ 含有把任意一对不同点 $r,r'$ 映到任意另一对这样的点、并以任意方式提升到 $\frL$ 纤维的元素。所以为了 $\G$-不变，$Z_{r,r'}$ 只能是 $\delta(r,r')$ 乘以 $\frL$ 纤维上的常值映射。这不是合理的答案。形式上，人们会把 $Z_{r,r'}$ 解释为
\be\label{nutto} Z_{r,r'}\stackrel{?}{=} \sum_{\alpha\in S} \bar\psi_\alpha(r) \psi_\alpha(r'), \ee
其中态 $\psi_\alpha$（$\alpha\in S$）构成 Hilbert 空间的一组基。要使这样的公式与 $Z_{r,r'}\sim\delta(r,r')$ 相容，就需要无限维的 Hilbert 空间，即使相空间 $M$ 是紧的。这不是量子理论应有的行为。而且对 $M=\R^2$（我们知道量子化应当意味着什么），并不存在像 \eqref{nutto} 那样的公式（记住 $r,r'$ 是相空间点，由位置与动量共同标记）。具体地说，若在微扰论中计算 $Z_{r,r'}$，首先发现的是作用量 $\sI$ 的经典 Euler--Lagrange 方程只告诉我们 $\dot x=0$，即映射 $x:I\to M$ 为常值。要计算 $Z_{r,r'}$ 中 delta 函数系数的一圈修正，需要把作用量在常值映射附近线性化所得的微分算符的行列式。但这个微分算符在“指定两端相空间点”的边界条件下不是自伴的，其行列式难以定义。

在 $\R^2$ 的例子中（坐标为 $p,q$，辛形式 $\omega=\d p\d q$），作用量为 $\sI=\int \d t\, p \dot q$。在解 $p=q=0$ 附近线性化，得到算符 $D=\begin{pmatrix} 0 & \frac{\d}{\d t}\cr -\frac{\d}{\d t}& 0\end{pmatrix}$。本征值问题 $D\psi=\lambda\psi$ 在“$p,q$ 在端点为零”的边界条件下没有非零解。问题在于这个边界条件并不使 $D$ 自伴。要得到自伴问题，每端只应指定一半边界值，更具体地说，所固定的变量应当 Poisson 交换。特别地，若 $\P$ 是第 \ref{gq} 节意义下 $M$ 的实极化，则可以要求端点值 $x(0)$ 与 $x(1)$ 分别落在 $\P$ 的指定叶 $F,F'$ 中。\footnote{\label{BKStwo}更一般地，可以取 $F$ 为一个极化 $\P$ 的叶，$F'$ 为另一个极化 $\P'$ 的叶。在这种情形下计算会得到 Blattner--Kostant--Sternberg 核，即脚注 \ref{BKS} 中提到的对象。}在此边界条件下计算路径积分，会回到几何量子化用极化 $\P$ 量子化的配方。所得结果本质上依赖于 $\P$。

尽管考察路径积分 \eqref{hiff} 并不能提供回避经典到量子过渡中反常的办法，这种量子化问题的表述仍然重要。假定我们有某种良好定义的微观理论，它以某种方式产生与 $M$ 相联系（也可能与进入该理论定义的其它数据相联系）的 Hilbert 空间 $\H$。再假定这个微观理论在区间 $I$ 上的路径积分可以约化为 \eqref{wiff} 加上一些可形式上论证为无关的“大质量”自由度。那么，在所讨论的微观理论确实良好定义的前提下，这实际上意味着：走向这个微观理论，给了我们一种把（\eqref{wiff} 的）病态量子化问题嵌入一个更大、良好定义的问题的办法。换言之，微观理论给了我们一种理解原病态问题的办法。因此，在这些条件下，把 $\H$ 定义为 $M$ 的量子化是合理的。我们讲到膜量子化时，实际采用的就是这个做法。

现在讨论\textbf{不}引入极化时能做什么。前文之所以考虑带边的一维流形（即单位区间 $I$）上的路径积分，是因为一般地，在 Feynman 路径积分中，初始或末态量子态出现在路径积分所在空间的边界上，这里就是 $I$。把 $I$ 换成圆周 $S^1$，就可以避免选取极化。此时考虑同样的路径积分，只是对环路 $x:S^1\to M$ 的空间 $\pP$ 积分。不需要边界条件，因为圆周没有边界。只要不引入边界，所计算的东西就不能与量子 Hilbert 空间中的态相联系。

在无法定义 Hilbert 空间的情况下，能期望计算什么？在形变量子化中——如 \cite{Many} 所提出，并在任意辛流形上严格实现 \cite{WL,Fedosov}——目标是把 $M$ 上函数构成的交换结合代数 $\A_0$ 形变为一个非交换但仍结合的代数 $\A$，形变按 $\hbar$ 的幂级数到所有阶。到 $\hbar$ 的一阶，形变应由 Poisson 括号给出。这里没有构造 $\A$ 在其上作用的 Hilbert 空间。\footnote{更精确地说，没有构造大小合适的 Hilbert 空间。$\A$ 自身可以用马上要引入的迹赋予 Hilbert 空间结构；即使 $M$ 紧，这个 Hilbert 空间也是无限维的。}

形变量子化未必与路径积分绑定，但圆周上的路径积分为形变量子化提供了自然的框架。形式框架如下（此框架已用 BV 量子化的机器严格化到 $\hbar$ 的所有阶 \cite{SL}）。设 $f$ 是 $M$ 上的函数，$r$ 是 $S^1$ 上的点。给定环路 $x:S^1\to M$，可把 $f$ 拉回到 $S^1$ 上的函数，并在点 $r$ 处取值，得到 $M$ 的环路空间 $\pP$ 上的函数，记作 $f(r)$。现在设 $f_1,f_2,\cdots, f_k$ 是 $M$ 上的函数，$r_1,r_2,\cdots, r_k$ 是 $S^1$ 上按循环次序排列的点，考虑路径积分
\be\label{cycl}\bigl\langle f_1 f_2 \cdots f_k\bigr\rangle =\int_\pP\D x \,e^{\i \sI} \,f_1(r_1) f_2(r_2)\cdots f_k(r_k). \ee
由于经典作用量在 $S^1$ 的保定向重参数化下不变，量 $\bigl\langle f_1 f_2 \cdots f_k\bigr\rangle$ 在形式上只应依赖点 $r_1,r_2,\dots,r_k$ 的循环次序，而不依赖具体点的选取。“关联函数” $\bigl\langle f_1 f_2 \cdots f_k\bigr\rangle$ 含有辛流形 $M$ 形变量子化所需的信息。关联函数与形变量子化的基本关系是：若
\be\label{woppo} \bigl\langle f_1 f_2 f_3\cdots f_k\bigr\rangle =\bigl\langle g f_3 \cdots f_k\bigr\rangle \ee
（其中 $g$ 依赖 $f_1,f_2$ 而不依赖 $f_3,\cdots,f_k$），则宣告在代数 $\A$ 中有 $\h f_1\h f_2=\h g$（其中 $\h f$ 是 $\A$ 中在 $\hbar\to0$ 极限下约化为 $f\in \A_0$ 的元素）。这个关系是算符乘积展开的一个版本，适用于拓扑场论。循环对称意味着关联函数 $\bigl\langle f_1 f_2 f_3\cdots f_k\bigr\rangle$ 具有迹的代数性质（例如 $\bigl\langle fg\bigr\rangle =\bigl\langle gf \bigr\rangle$），尽管它实际上并不是 $\A$-模上的迹。

至少就目前的理解而言，形变量子化破坏或至少形变了经典力学的对称性，其方式与量子化有些相似，尽管细节不同。在 Fedosov 的进路 \cite{Fedosov} 中，量子形变代数 $\A$ 的构造依赖于流形 $M$ 上辛联络的选取（即 $M$ 的切丛上使辛形式协变常数的联络）。每个 $M$ 上都有辛联络（不像极化 \cite{Gotay}），但辛联络远不唯一。一般的辛同构会把用一个辛联络 $\Gamma$ 构造的代数 $\A_\Gamma$ 映到用另一个辛联络 $\Gamma'$ 构造的代数 $\A_{\Gamma'}$。$\A_\Gamma$ 与 $\A_{\Gamma'}$ 同构，但并非典则同构（见 \cite{Fedosov} 推论 4.5）。不同的可能 $\A$ 之间没有典则同构，意味着 $M$ 的辛同构群不作用在任何一个 $\A$ 上。若 $G$ 是任一\textbf{紧}辛同构群，则存在 $G$-不变的辛联络，因此任何这样的 $G$ 都可实现为量子代数 $\A$ 的对称群。相比之下，几何量子化一般不能定义成在紧辛同构群下不变。\footnote{例如，若 $M$ 是带平移不变辛形式 $\omega$ 的二维环面，则可在 $M$ 上取平移不变的（平坦）辛联络。因此形变量子化可以定义成保持环面平移不变性的。但环面的几何量子化需要选取带联络的预量子线丛 $\frL$，而任何选取都破坏平移对称性。这是因为 $\frL$ 的曲率使 $\frL$ 绕 $M$ 中不可缩闭路的和乐不可能平移不变。}

与量子化不同，形变量子化只依赖辛结构 $\omega$，而不依赖预量子线丛 $\frL$ 的存在或选取。事实上，就形变量子化而言，作用量 $\sI=\oint_{S^1}x^*(\sT )/\hbar$ 可以改写成只涉及 $\omega$ 而不涉及 $\sT$ 的形式。只需注意，在 $\hbar$ 的微扰论中可假定 $x:S^1\to M$ 几乎为常值。设 $D$ 是以 $S^1$ 为边界的圆盘。几乎常值的映射 $x:S^1\to M$ 可以通过把 $x$ 的像收缩为 $M$ 中一点而延拓为映射 $x:D\to M$。此延拓不唯一，但在同伦意义下唯一。于是可以把 $\frac{1}{\hbar}\oint_{S_1}x^*(\sT )$ 换成 $\int_D  x^*(\omega)$，并把 \eqref{cycl} 改写成
\be\label{pycl}\bigl\langle f_1 f_2 \cdots f_k\bigr\rangle =\int_\pP\D x\, \exp\left(\i \int_D  x^*(\omega)\right) f_1(r_1) f_2(r_2)\cdots f_k(r_k), \ee
它显然只依赖 $\omega$ 而不依赖 $\frL$。

形变量子化在更一般的、未必辛的 Poisson 流形语境中也有意义。Poisson 流形按定义有由公式 $\{f,g\}=\sum_{ij}r^{ij}\partial_i f \partial_j g$ 给出的 Poisson 括号，其中 $r^{ij}$ 是反对称张量，满足通常的 Jacobi 恒等式，但 $r^{ij}$ 未必可逆（在辛流形上 $r^{ij}$ 是辛形式的逆，特别地可逆）。在这一语境中仍可给出形变量子化的路径积分描述，但必须使用二维而非一维路径积分 \cite{Kon,CF,CF2}。
